%% Finite splitting, back-and-forth, and the Cantor-space point map.

\section{The countable atomless Boolean algebra}

\begin{annexitem}{annex:p10.atomless.1}{thm:p10-atomless-extension}{Extending finite Boolean isomorphisms}
\begin{proof}
\paragraph{Finite Boolean algebras and their atoms.}
Let \(F\) be finite and nontrivial.
For every \(b\ne0\), the finite set of nonzero elements below \(b\)
has a minimal member, which is an atom.
Distinct atoms are disjoint: a nonzero meet below both would
force them to be equal.
Every \(b\in F\) is the join of the atoms below it.
Indeed, if their join \(d\) were strictly below \(b\), then
\(b\wedge\neg d\ne0\) would contain a further atom below \(b\).
Thus the atoms form a partition of \(1\), and different subsets
of the atom set have different joins.
Taking joins identifies \(F\) with the power-set Boolean algebra
of its atoms.
In particular, a bijection of finite atom partitions extends
uniquely to a Boolean isomorphism.
This also proves that no finite nontrivial Boolean algebra is atomless.

A subalgebra generated by \(m\) elements \(b_1,\ldots,b_m\) is
finite.  To see this, form the \(2^m\) Boolean cells
\[
 \bigwedge_{j=1}^m b_j^{\varepsilon_j},
 \qquad
 b_j^1=b_j,\quad b_j^0=\neg b_j,\quad
 \varepsilon_j\in\{0,1\}.
\]
The nonzero cells are disjoint and have join \(1\).
Their joins form a subalgebra containing each generator, while
each cell is a Boolean term in the generators.
They therefore generate exactly that finite subalgebra, whose
size is at most \(2^{2^m}\).
For \(m=0\) the sole cell is the empty meet \(1\), and the
generated subalgebra is \(\{0,1\}\).

\paragraph{Splitting a nonzero element.}
In an atomless algebra \(C\), any nonzero \(c\) can be partitioned
into \(k\) nonzero pieces for every positive integer \(k\).
For \(k=1\), use \(c\).
To increase the number of pieces, split a nonzero piece \(d\)
using \(0<e<d\):
\[
 d=e\vee(d\wedge\neg e).
\]
Both parts are nonzero and disjoint.
Replacing one piece by these two proves the assertion by induction.

\paragraph{Extending the atom map.}
Let \(p_1,\ldots,p_r\) be the atoms of \(F\).
Each is the join of precisely the atoms of \(F'\) below it.
There is at least one such atom; enumerate them as
\(p_{i,1},\ldots,p_{i,k_i}\).
Every atom of \(F'\) lies below exactly one \(p_i\), since
the \(p_i\) partition \(1\).
The isomorphism \(h\) sends the \(p_i\) to the atoms \(h(p_i)\)
of \(G\).
For each \(i\), partition \(h(p_i)\) inside \(C\) into
\(k_i\) nonzero pieces \(c_{i,1},\ldots,c_{i,k_i}\).
Together these pieces form a finite partition of \(1_C\).

Let \(G'\) be the subalgebra consisting of all joins of these pieces,
and define
\[
 h'\left(\bigvee_{(i,j)\in E}p_{i,j}\right)
       =\bigvee_{(i,j)\in E}c_{i,j}.
\]
Unique decomposition into atoms makes this a well-defined
isomorphism \(F'\to G'\).
It maps \(p_i\) to \(\bigvee_jc_{i,j}=h(p_i)\), so it extends
\(h\) on every join of atoms of \(F\), namely on all of \(F\).
These equalities also show \(G\leq G'\).

For a one-element extension
\(F'=\langle F\cup\{b\}\rangle\), the pieces of each \(p_i\)
are the nonzero members of
\(\{p_i\wedge b,p_i\wedge\neg b\}\).
There are one or two of them, and the construction makes
exactly the same split on the corresponding target atom.
\end{proof}
\end{annexitem}

\begin{annexitem}{annex:p10.atomless.2}{thm:p10-countable-atomless-uniqueness}{Back-and-forth uniqueness}
\begin{proof}
By Annex~\ref{annex:p10.atomless.1}, nontrivial atomless \(B,C\)
cannot be finite.
Choose enumerations
\[
 B=\{b_0,b_1,\ldots\},\qquad C=\{c_0,c_1,\ldots\}.
\]
Start with the prescribed isomorphism \(h_0:F_0\to G_0\)
between finite subalgebras.
For the unqualified uniqueness statement use
\(F_0=\{0_B,1_B\}\), \(G_0=\{0_C,1_C\}\) and the
isomorphism matching the constants.

Construct increasing finite subalgebras and extending isomorphisms
\(h_n:F_n\to G_n\).
At the forth step, the generated subalgebra
\[
 F'=\langle F_{2n}\cup\{b_n\}\rangle
\]
is finite.  The extension theorem, with atomless target \(C\),
extends \(h_{2n}\) to
\(h_{2n+1}:F'\to G_{2n+1}\).
Set \(F_{2n+1}=F'\).
At the back step, form
\[
 G'=\langle G_{2n+1}\cup\{c_n\}\rangle.
\]
Apply the extension theorem to \(h_{2n+1}^{-1}\), with
atomless target \(B\), extending it to an isomorphism
\(G'\to F_{2n+2}\).
Its inverse is \(h_{2n+2}:F_{2n+2}\to G_{2n+2}\), where
\(G_{2n+2}=G'\).
This inverse extends \(h_{2n+1}\).

Every element of \(B\) eventually enters a domain and every
element of \(C\) eventually enters a range.
Consequently \(h=\bigcup_n h_n\) is a total bijection \(B\to C\).
Any finitely many arguments of a Boolean operation occur in
one \(F_n\), where \(h_n\) preserves that operation.
The constants already occur in \(F_0\).
Thus \(h\) is a Boolean isomorphism extending \(h_0\).

The construction can be made with no sequence of arbitrary choices:
use the fixed enumerations to order each finite atom set, and
always split the first current piece in enumeration order.
For the split, choose the first enumerated element strictly
between \(0\) and that piece.
Atomlessness guarantees such an element.
This specifies the successive finite extensions recursively.
\end{proof}
\end{annexitem}

\begin{annexitem}{annex:p10.atomless.3}{thm:p10-cantor-cylinder-algebra}{Existence through clopen cylinders}
\begin{proof}
The all-zero sequence proves \(K\ne\varnothing\).
The cylinders form a basis for the product topology: any restriction
on finitely many coordinates can be refined by fixing an initial
segment containing them all.
For each \(n\), the \(2^n\) cylinders of length \(n\) are a
partition of \(K\) into nonempty sets.
The complement of one cylinder is the union of the other
cylinders of that length, so each cylinder is clopen.
Two different sequences differ at a coordinate; their initial
segments through that coordinate give disjoint neighborhoods.
Thus \(K\) is Hausdorff.

\paragraph{Compactness by a binary-tree argument.}
Let \(\mathcal U\) be an open cover and suppose it has no finite
subcover.
Call a finite word \(s\) bad if \([s]\) cannot be covered by
finitely many members of \(\mathcal U\).
The empty word is bad.
Since \([s]=[s0]\cup[s1]\), a bad word has a bad child:
otherwise two finite covers would combine to cover \([s]\).
Starting at the empty word, choose child \(0\) if it is bad
and child \(1\) otherwise.
This recursion defines \(x\in K\) all of whose initial segments
are bad.
Some \(U\in\mathcal U\) contains \(x\); openness supplies a
cylinder \([x|_n]\subseteq U\).
That cylinder has a one-member cover from \(\mathcal U\),
a contradiction.
Hence \(K\) is compact.

\paragraph{All clopens occur at a finite level.}
Let \(V\subseteq K\) be clopen.
Take all cylinders contained in \(V\); they cover \(V\).
Together with the open set \(K\setminus V\), they cover \(K\).
Compactness gives finitely many of those cylinders covering \(V\),
so \(V\) is their union.
The empty clopen is the empty union.
Refine every cylinder in a nonempty finite union to a common
length \(n\), using
\[
 [s]=\bigcup\{[t]:t\in\{0,1\}^n,\ s\text{ is a prefix of }t\}.
\]
Thus \(V\in\mathcal B_n\).
Conversely every member of \(\mathcal B_n\) is clopen, proving
\(\operatorname{Clop}(K)=\mathcal B_{\rm cyl}\).

Because the length-\(n\) cylinders are disjoint and nonempty,
their unions form a Boolean algebra of size \(2^{2^n}\),
with exactly those cylinders as its atoms.
Each \([s]\) is the disjoint union of the two nonempty
cylinders \([s0],[s1]\).
Therefore
\(\mathcal B_n\subsetneq\mathcal B_{n+1}\).
Their union is countable: a member has a description by an
integer \(n\) and a subset of the explicitly finite set
\(\{0,1\}^n\), and these descriptions can be enumerated.
The strict inclusions show the union is infinite.
It is nontrivial since \(\varnothing\ne K\).

If \(V\in\mathcal B_{\rm cyl}\) is nonempty, choose a cylinder
\([s]\) in a finite-level representation of \(V\).
Then \([s0]\) is a nonempty proper subset of \(V\):
the disjoint nonempty set \([s1]\) is still contained in \(V\).
This gives \(0<[s0]<V\) in the Boolean order, proving atomlessness.

For completeness, the direct-limit assertion has its usual
universal meaning.  Let \(D\) be a Boolean algebra and suppose homomorphisms
\(g_n:\mathcal B_n\to D\) satisfy
\(g_{n+1}|_{\mathcal B_n}=g_n\).
Their union is a well-defined homomorphism
\(\mathcal B_{\rm cyl}\to D\), because every finite tuple
lies in some common level.
It is the unique map extending all \(g_n\).
This proves the direct-limit property of the inclusion sequence.
The uniqueness theorem now identifies every nontrivial countable
atomless Boolean algebra with this example.
\end{proof}
\end{annexitem}

\begin{annexitem}{annex:p10.atomless.4}{thm:p10-countable-atomless-stone}{Identifying the Stone space with Cantor space}
\begin{proof}
Existence and uniqueness of the algebra provide
\(\alpha:B\to\mathcal B_{\rm cyl}\).
For each \(x\in K\), membership at \(x\) preserves Boolean
operations from \(\mathcal B_{\rm cyl}\) to \(\{0,1\}\).
Its inverse image of \(1\), transported by \(\alpha\), is the
proper ultrafilter \(H(x)\).
Thus the displayed point map is well-defined.

Let \(U\) be any ultrafilter of \(B\), and transfer it to the
ultrafilter
\[
 V=\{D\in\mathcal B_{\rm cyl}:\alpha^{-1}(D)\in U\}.
\]
An ultrafilter contains exactly one member of each finite partition
of \(1\) into nonzero Boolean elements.
It contains at most one because distinct pieces have meet \(0\).
If it contained none, it would contain every complement, whose
finite meet is \(0\), contradicting properness.
For each \(n\), therefore, there is a unique word \(s_n\) of
length \(n\) with \([s_n]\in V\).
The selected cylinders at consecutive levels intersect, so
\(s_{n+1}\) extends \(s_n\).
They define a unique sequence \(x\in K\).

For any \(D\in\mathcal B_{\rm cyl}\), choose a level \(n\)
at which it is a union of cylinders.
Then
\[
 D\in V
 \quad\Longleftrightarrow\quad
 [s_n]\subseteq D
 \quad\Longleftrightarrow\quad x\in D.
\]
For the first implication, \(D\) and \([s_n]\) in \(V\)
have nonempty intersection, and a level-\(n\) cylinder is either
contained in \(D\) or disjoint from it.
The converse follows from upward closure of \(V\).
Thus \(U=H(x)\), proving surjectivity.
If \(x\ne y\), a cylinder contains exactly one of them, and
its preimage under \(\alpha\) belongs to exactly one of
\(H(x),H(y)\).
Hence \(H\) is injective.

For each \(b\in B\),
\[
 H^{-1}(\widehat b)=\{x:x\in\alpha(b)\}=\alpha(b),
\]
which is open, so \(H\) is continuous by the definition of the
Stone topology.
For every cylinder \([s]\), bijectivity gives
\[
 H([s])=\widehat{\alpha^{-1}([s])},
\]
which is open in \(X_B\).
Since cylinders form a basis, \(H\) is open and its inverse
is continuous.
This proves the homeomorphism directly.
\end{proof}
\end{annexitem}

\begin{annexitem}{annex:p10.atomless.5}{thm:p10-atoms-isolated}{Atoms and isolated points}
\begin{proof}
Let \(p\) be an atom.
The set \(U_p=\{b:p\leq b\}\) is a proper filter.
For any \(b\), the element \(p\wedge b\) is either \(0\)
or \(p\).
In the first case \(p\leq\neg b\), and in the second
\(p\leq b\); both cannot hold.
Thus \(U_p\) contains exactly one of \(b,\neg b\), and is
an ultrafilter.
Any ultrafilter containing \(p\) contains every \(b\geq p\);
if \(p\not\leq b\), then \(p\wedge b=0\), so it cannot
contain \(b\).
It is therefore \(U_p\), and
\(\widehat p=\{U_p\}\) is clopen.

Conversely, let \(U\) be isolated in \(X_B\).
The sets \(\widehat b\) form a basis, since finite intersections
are represented by meets.
Choose \(b\) with
\[
 U\in\widehat b\subseteq\{U\}.
\]
Here \(b\ne0\), since no proper ultrafilter contains \(0\).
If \(b\) were not an atom, choose \(0<c<b\).
Then \(d=b\wedge\neg c\) is also nonzero.
By the Boolean prime ideal theorem the proper filters generated
by \(c\) and \(d\) extend to ultrafilters \(U_c,U_d\).
Both contain \(b\), while one contains \(c\) and the other
contains \(\neg c\), so they are distinct.
This contradicts \(\widehat b=\{U\}\).
Thus \(b\) is an atom and \(U=U_b\).
Distinct atoms give distinct principal ultrafilters, since their
meet is \(0\); this proves the claimed bijection.

If \(B\) is finite and nontrivial, its atoms form a finite partition
of \(1\), so each ultrafilter contains one atom and is the
corresponding principal ultrafilter.
All its Stone points are consequently isolated.
If there are \(k\) atoms, the finite-atom decomposition proves
\(|B|=2^k\), and the Stone space is discrete with \(k\) points.
For the one-element algebra there are no proper ultrafilters and
no atoms; its Stone space is empty.

Finally every cylinder \([s]\) in \(K\) contains \([s0]\)
and \([s1]\), both nonempty.
No point of \(K\) can therefore be an open singleton.
This agrees with the atomlessness of its clopen algebra and with
the explicit homeomorphism proved above.
\end{proof}
\end{annexitem}
